\documentclass[12pt]{scrartcl}

%Matthis Vorlage
\usepackage{../mathe_definition}
%wird in Matthis Volage auch gemacht aber das rafft mein texstudio nicht und highlighted Dinge falsch
\usepackage{amsmath}

\usepackage[shortlabels]{enumitem}

%Ausgelagerter Code fürs schöne seitenlayout
\date{\today}
\newcommand{\Blattnummer}{%
	6%
}

\title{{\itshape Stochastische Prozesse}}
\module{Blatt 1}
\author{Philip Wunderlich \and Jonathan Langer}
%\matriculation{3737927\and 3750030}

\fancypagestyle{netter_header}{
	\renewcommand{\and}{\\}
	\fancyhead[L]{\MyAuthor}
	\fancyhead[C]{\MyTitle\\ \MyModule}
	\fancyhead[R]{\MyDate\\Seite \thepage{} von \pageref{LastPage}}
	\fancyfoot[C]{}
}
\pagestyle{netter_header}

%Abkürzungen
\newcommand{\einruecken}{\begin{addmargin}[25pt]{0pt}} 
	%LaTeX mag keine Umlaute :(
\newcommand{\ausruecken}{\end{addmargin}}
\newcommand{\R}{\ensuremath{\mathbb{R}}}
\newcommand{\Om}{\ensuremath{\Omega}}
\newcommand{\om}{\ensuremath{\omega}}
\newcommand{\A}{\ensuremath{\mathcal{A}}}
\newcommand{\B}{\ensuremath{\mathcal{B}}}
\newcommand{\C}{\ensuremath{\mathcal{C}}}
\newcommand{\WraumAnn}{\ensuremath{(\Om, \A, P)} ein Wahrscheinlichkeitsraum}
\newcommand{\MraumAnn}{\ensuremath{(E,\B)} ein Messraum}
\newcommand{\Xt}{\ensuremath{(X_t)_{t\in T}}}
\newcommand{\Yt}{\ensuremath{(Y_t)_{t\in T}}}
%\newcommand{\Xw}{\ensuremath{X_t(\omega)}}
%\newcommand{\Yw}{\ensuremath{Y_t(\omega)}}
%\newcommand{\Xq}{\ensuremath{X_q(\omega)}}
%\newcommand{\Yq}{\ensuremath{Y_q(\omega)}}
\newcommand{\E}{\ensuremath{\mathbb{E}}}
\newcommand{\1}{\ensuremath{\textbf{1}}}
\newcommand{\LwieLp}{\ensuremath{\mathcal{L}}}
% \L gibts anscheinend schon
\newcommand{\scriptL}{\ensuremath{\mathcal{L}}}
\newcommand{\scriptP}{\ensuremath{\mathcal{P}}}
\newcommand*\diff{\mathop{}\!\mathrm{d}}
\newcommand{\F}{\ensuremath{\mathcal{F}}}
\newcommand{\Dn}{\ensuremath{D_n}}
\newcommand{\cn}{\ensuremath{c_n}}
\newcommand{\ep}{\ensuremath{{e^{tc_n}}}}
\newcommand{\en}{\ensuremath{{e^{-tc_n}}}}
\newcommand{\N}{\ensuremath{\mathbb{N}}}
%\newcommand{\Xn}{\ensuremath{(X_n)_{n\in \N}}}
%\newcommand{\Yn}{\ensuremath{(Y_n)_{n\in \N}}}
\newcommand{\var}{\text{var}}
\newcommand{\textP}{\text{P}}
\newcommand{\K}{\ensuremath{\mathcal{K}}}

\begin{document}

\subsection*{Excercise 19}
\subsubsection*{Lemma 2.5.2}

\begin{enumerate}[$i)$]
	%Setzt Counter schon zu Beginn auf 1, damit danach das erste Item als ii) kommt.
	\setcounter{enumi}{1}
	\item For $\om\in\Om$ define 
	$$U_{a,b}^N(\om):=\sup\{k\in\N:\sigma_k(\om)\le N\}$$
	with $\sigma_k$ stopping times for all $k\ge0$ and $\sigma_k\le\sigma_l$ for $k\ge l$. The resulting map $U_{a,b}^N: \Om\rightarrow\overline{\N}_0$ is measurable.
\end{enumerate}

\subsubsection*{Proof}
\einruecken

trivial

\ausruecken

\subsection*{Excercise 20}
\subsubsection*{Lemma 2.6.4}

Let $(\Om,\A,P)$ be a probability space and $\K\subset\scriptL_0(\A)$ be nonempty. If there exists $Y\in\scriptL_1(P)$ with
$$|X|\le Y,\quad X\in\K,$$
then $\K$ ist uniformly integrable.

\subsubsection*{Proof}
\einruecken

By MWT lecture, we know that $N:=\{Y=\infty\}$ has measure $0$. For $X\in\K$ and $\om\in\Om\setminus N$, we have $$\om\in\{|X|<|X(\om)+1|\},\quad |X(\om)|\le Y(\om)<\infty.$$
Therefore, $\{|X|<c\}\cap(\Om\setminus N)\nearrow\Om\setminus N$ and
$$Y\1_{\{|X|<c\}\cap(\Om\setminus N)}\nearrow Y\1_{\Om\setminus N}.$$
Because integrals don't change on nullsets, this yields
$$\lim_{c\rightarrow\infty}\int_{|X|<c} Y\diff P
= \lim_{c\rightarrow\infty}\int_{\Om} Y\1_{\{|X|<c\}\cap(\Om\setminus N)}\diff P
= \int_{\Om}Y\1_{\Om\setminus N}\diff P
= \int_{\Om}Y\diff P$$
with the middle equality using Beppo-Levi. Further, the equality
$$\int_{\Om}Y\diff P =\int_{|X|<c}Y\diff P +\int_{|X|\ge c}Y\diff P$$
now yields
$$\lim_{c\rightarrow\infty}\int_{|X|\ge c}Y\diff P =0$$
by rearranging and using the definition of the limit. Therefore, we finally get
\begin{align*}
\lim_{c\rightarrow\infty}\sup_{X\in\K}\int_{|X|\le c}|X|\diff P
&\
\end{align*}

....hmmm iwi funktioniert das actually gar nicht..

\ausruecken

\subsubsection*{Theorem 2.6.5}

Let $(\Om,\A,P)$ be a probability space and $\K\subset\scriptL_0(\A)$ be nonempty. Further, let $\phi:[0,\infty)\rightarrow[0,\infty)$ be measurable with
$$\lim_{t\rightarrow\infty}\frac{\varphi(t)}{t}=\infty.$$
If
$$B:=\sup_{X\in\K}\E\varphi(|X|)<\infty,$$
then $\K$ ist uniformly integrable.

\subsubsection*{Proof}
\einruecken

trivial

\ausruecken

\subsection*{Excercise 21}
\subsubsection*{Corollary 2.6.9}

Let $(\Om,\A,P)$ be a probability space, $(X_n)_{n\ge 1}\subset\scriptL_1(P)$ and $X\in\scriptL_0(\A)$ with $X_n\rightarrow X$ in probability. If there exists $Y\in\scriptL_1(P)$ with $|X_n|\le Y$ for all $n\ge 1$, then
$$||X_n-X||_{\scriptL_1(P)}\rightarrow 0.$$
In particular, $\E X_n\rightarrow\E X$.

\subsubsection*{Proof}
\einruecken

By \textit{Lemma 2.6.4}, proved in Exercise 20, $\{X_n:n\ge 1\}$ is uniformly integrable.
\textit{Theorem 2.6.8.} now directly yields the convergence $X_n\rightarrow X$ in $\scriptL_1(P)$.
The proposition $\E X_n\rightarrow\E X$ follows from
$$|\E X_n-\E X|=|\E(X_n-X)|\le \E|X_n-X|=||X_n-X||_{\scriptL_1(P)}.$$

\ausruecken

\subsubsection*{Corollary 2.6.10}

Let $(\Om,\A,P)$ be a probability space, $(X_n)_{n\ge 1}\subset\scriptL_p(P)$ with $p\in(1,\infty]$ and $X\in\scriptL_0(\A)$ with $X_n\rightarrow X$ in probability. If
$$\sup_{n\ge 1}||X_n||_{\scriptL_p(P)}<\infty,$$
then $\E X_n\rightarrow\E X$.

\subsubsection*{Proof}
\einruecken

Using \textit{Theorem 2.6.5} from Exercise 20 with $\varphi(t):=t^p$, we get the uniformly integrability of $\{X_n:n\ge 1\}$. Again, we get $\scriptL_1(P)$-convergence and $\E X_n\rightarrow\E X$ by using \textit{Theorem 2.6.8}.

\ausruecken
	
\end{document}